\setlength{\parindent}{0pt}
% ------------------------------------------------------------
\chapter{Introduction}
\label{chap:intro}
% ------------------------------------------------------------
\section{Side-Channel Attacks and the Masking Countermeasure}
\label{sec:intro-masking}

A cryptographic scheme can be secure on paper and insecure on the bench. When an algorithm
runs on a real device its power draw, its timing, and its electromagnetic
emanations all vary with the data it handles, and an adversary who measures those
physical quantities can recover a key the mathematics was supposed to protect.
This is called a \emph{side-channel attack}, and differential power analysis
\parencite{kocher1999differential} is one type of such an attack. By correlating many power
traces against a guess of an intermediate value, it extracts secret keys from
implementations whose black-box behaviour is beyond reproach. The insecurity lies not in
the cipher but in the way the hardware computes it.

\emph{Masking} is the standard algorithmic countermeasure
\parencite{chari1999towards} against differential power analysis. In general it
splits each secret value into $n$ shares so that any $n-1$ of them are independent
of the secret, and carries the computation out on the shares rather than on the
secret itself. The shares are combined by an operation fixed in advance, and the
choice of operation depends on the masking scheme. Arithmetic masking adds the shares in a ring
and multiplicative masking multiplies them, but the form we adopt is \emph{Boolean
masking}, where the operation is exclusive-or over $\mathbb{F}_2$. In a masked
circuit, a secret $x$ is then written as $x = x_1 \oplus \dots \oplus x_n$ with the
first $n-1$ shares drawn uniformly at random, so that any $n-1$ of them together
reveal nothing about $x$. The computation runs on the shares, and because any
$n-1$ of them are independent of the secret, physical leakage that reflects only
some of the shares tells an adversary nothing on its own. To learn anything, the adversary
must combine leakage from several shares at once, and each extra share raises the
measurement cost sharply. Realising this
masking across a whole cipher is delicate, since the non-linear operations of the (unmasked)
algorithm force the shares of different variables to interact in the masked circuit, and it is precisely
those interactions that a masking scheme must handle with care \parencite{10.1007/978-3-540-45146-4_27, rivain2010provably}.

\section{Verifying Masked Implementations}
\label{sec:intro-verification}

% Masking is easy to get subtly wrong. Splitting a secret into shares helps only as
% long as the shares are never brought back together somewhere an attacker can
% reach, and a real implementation gives them many chances to. A single wire may
% recombine two shares of the same secret, or a random mask meant to hide one value
% may cancel against a copy of itself elsewhere, and either slip reopens the very
% leak the masking was meant to close. There is a further subtlety. A masking must
% hold up against an attacker who combines the leakage of several intermediate
% values at once, not just one, so its correctness cannot be judged value by value;
% it depends on how the values leak jointly. The number of such combinations grows
% quickly with the size of the circuit and the number of shares, which puts checking
% by hand out of reach for all but the smallest implementations.

% What one wants, then, is a tool that takes a masked circuit and decides,
% mechanically and for all of those combinations at once, whether any of them leaks.
% This is the problem of \emph{formally verifying} a masked implementation, and a
% line of work has grown around it \parencite{barthe2015verified,
% 10.1007/978-3-030-29959-0_15}. The condition to establish is one about
% distributions, that the values on every small set of wires are distributed the
% same way whatever the secret is, so that nothing about the secret can be recovered
% from them. A verifier must certify this for all such sets together, and two
% difficulties sit underneath. Even one set is already a statement about a
% distribution over all the internal random choices, and there are combinatorially
% many sets to settle. The tools that scale do not compute those distributions
% head-on. They rewrite each set of values into a form in which the secrets have
% visibly cancelled, and they organise the search so that a single such certificate
% settles many sets at once. This thesis builds a verifier in that constructive
% spirit.

\section{Scope and Threat Model}
\label{sec:intro-scope}

We work in the \emph{probing model} of Ishai, Sahai, and Wagner
\parencite{10.1007/978-3-540-45146-4_27}. The circuit is arithmetic over
$\mathbb{F}_2$, built from two-input XOR and AND gates, and its inputs are
partitioned into secret, random, and public wires. An adversary of order $d$
chooses any $d$ wires and learns their joint values over the draw of the internal
randomness, and the circuit is \emph{$d$-probing secure} when no such choice tells
the secrets apart (Chapter~\ref{chap:preliminaries}). The leakage is
\emph{value-based}: a probed wire leaks the logical value it carries and nothing
more. Physical effects that make a wire leak about more than its final value, such
as glitches and transitions in hardware \parencite{10.1007/978-3-030-29959-0_15},
are outside our model, as is any question of composing separately verified gadgets
into a larger one. The object of study is the exact probing security of a single
circuit, decided over its wire values.

Within that model our aim is exact verification rather than a sound
approximation. When the tool reports a circuit secure, every one of its
$\sum_{i \le d} \binom{|\mathcal{W}|}{i}$ observation sets has been accounted for,
and when it reports a circuit insecure it returns a specific leaking observation
as a witness, up to the one incompleteness we flag in
Chapter~\ref{chap:design}. The decision problem underlying a single observation is
hard in general, in a sense Chapter~\ref{chap:couplings} makes precise, which is
what forces the constructive route the tool takes rather than a direct count.

\section{Contributions}
\label{sec:intro-contributions}

This thesis presents \toolname, a verifier of $d$-probing security for masked
arithmetic circuits over $\mathbb{F}_2$, implemented in Lean~4. Its design is
organised around three components, and the contributions follow that division.

\textbf{A constructive oracle for one observation.} We give a rewrite engine that
decides secret-independence of a single probe by substituting away its randoms one
at a time, driving the observation toward a secret-free form
(Chapter~\ref{chap:design}). The engine is layered, a cheap single-occurrence rule
first and a complete substitution loop behind it, and we prove its verdicts sound.
To reach randoms that ordinary rewriting cannot, we add a \emph{multiplicative
factoring} transformation that reshapes products to expose masks buried inside
them, and we show it strictly enlarges the class of circuits the engine can
certify.

\textbf{An exhaustive traversal of the observation space.} We give a driver that
certifies every $d$-subset of the wires without enumerating them one by one,
splitting the space recursively in the manner of \texttt{maskVerif}
\parencite{10.1007/978-3-030-29959-0_15} and proving the split exhaustive through
a set-level reading of Vandermonde's identity (Chapter~\ref{chap:design}). One
discharge of a large union can prune a whole branch at once.

\textbf{A coupling view of rewriting, and probe-set extension.} We read a
successful discharge not merely as a yes-or-no verdict but as a
measure-preserving \emph{coupling} of the randomness, in the tradition of coupling
proofs for probabilistic programs \parencite{barthe2017coupling}, developed for
probing security in Chapter~\ref{chap:couplings}. That reading lets one
certificate blind many wires at once. We use it to grow a certified probe into a
whole \emph{safe set} of wires, so a single discharge settles many observations,
and we compare three mechanisms for the growth, adopting the coupling-based one and
proving the other two sit inside it (Chapter~\ref{chap:design}).

\textbf{Implementation and evaluation.} We describe the data structures that make
the engine efficient, a hash-consed circuit representation and a bit-sliced,
staged test for the extension (Chapter~\ref{chap:implementation}), and we evaluate
\toolname\ on a suite of masked gadgets, reporting its verification outcomes and
analysing where the search does redundant work (Chapter~\ref{chap:evaluation}).
Along the way the tool rediscovers a genuine flaw in a four-share gadget, which we
use as a case study.

\section{Thesis Outline}
\label{sec:intro-outline}

Chapter~\ref{chap:preliminaries} fixes notation and defines arithmetic circuits,
the probing model, and $d$-probing security, and studies the randomness
substitution technique we use througout. Chapter~\ref{chap:related} surveys the
verification literature we build on and depart from. Chapter~\ref{chap:couplings}
develops the coupling view of probing security and situates the underlying
decision problem in the complexity hierarchy. Chapter~\ref{chap:design} is the
core of the thesis, presenting our algorithm with its three components together with their algorithms and
correctness proofs, while Chapter~\ref{chap:implementation} gives the concrete data
structures and implementation details behind them. Chapter~\ref{chap:evaluation} reports the experiments.
Chapter~\ref{chap:futurework} outlines the directions our approach does not treat, and
Chapter~\ref{chap:conclusion} concludes the thesis.
